\ifdefined\gkver\else\def\gkver{student}\fi
\documentclass[\gkver]{gaokaozhenti}
\begin{document}

\chapter{2023年海南}

\begin{enumerate}
\item
%% number: 1
%% paperName: 2023年海南高考第 1 题 3 分
%% typeId: 1
%% type: 单选题
%% chapter: 原子和原子核
%% point: 天然放射性
%% method: 无
%% score: 3
%% degree: 950
%% duplicateId: 0
%% body:
钍元素衰变时会放出 $\beta$ 粒子，其中 $\beta$ 粒子是\xzanswer{C}

\fourchoices[answer=C]
{中子}
{质子}
{电子}
{光子}

%% answer: C
%% memo:
\memoanswer{
	放射性元素衰变时放出的三种射线 $\alpha$、$\beta$、$\gamma$ 分别是氦核流、电子流和光子流。

	故选 C。
}

\item
%% number: 2
%% paperName: 2023年海南高考第 2 题 3 分
%% typeId: 1
%% type: 单选题
%% chapter: 磁场
%% point: 洛伦兹力
%% method: 无
%% score: 3
%% degree: 850
%% duplicateId: 0
%% body:
如图所示，带正电的小球竖直向下射入垂直纸面向里的匀强磁场，关于小球运动和受力说法正确的是\xzanswer{A}

\onepicture[w=4.50cm]{figs/02.png}

\fourchoices[answer=A]
{小球刚进入磁场时受到的洛伦兹力水平向右}
{小球运动过程中的速度不变}
{小球运动过程的加速度保持不变}
{小球受到的洛伦兹力对小球做正功}

%% answer: A
%% memo:
\memoanswer{
	A．根据左手定则，可知小球刚进入磁场时受到的洛伦兹力水平向右，A 正确；

	BC．小球受洛伦兹力和重力的作用，则小球运动过程中速度、加速度大小，方向都在变，BC 错误；

	D．洛伦兹力永不做功，D 错误。

	故选 A。
}

\item
%% number: 3
%% paperName: 2023年海南高考第 3 题 3 分
%% typeId: 1
%% type: 单选题
%% chapter: 力物体的平衡
%% point: 动态平衡
%% method: 无
%% score: 3
%% degree: 750
%% duplicateId: 0
%% body:
如图所示，工人利用滑轮组将重物缓慢提起，下列说法正确的是\xzanswer{B}

\onepicture[w=5.00cm]{figs/03.png}

\fourchoices[answer=B]
{工人受到的重力和支持力是一对平衡力}
{工人对绳的拉力和绳对工人的拉力是一对作用力与反作用力}
{重物缓慢拉起过程，绳子拉力变小}
{重物缓慢拉起过程，绳子拉力不变}

%% answer: B
%% memo:
\memoanswer{
	AB．对人受力分析如图，

	\onepicture[w=3.50cm]{figs/03_an1.png}

	则有 $F_{N} + F_{T} = mg$

	其中工人对绳的拉力和绳对工人的拉力是一对作用力与反作用力，A 错误、B 正确；

	\onepicture[w=3.50cm]{figs/03_an2.png}

	CD．对滑轮做受力分析如图，则有 $F_{T} = \frac{mg}{2\cos\theta}$

	则随着重物缓慢拉起过程，$\theta$ 逐渐增大，则 $F_{T}$ 逐渐增大，CD 错误。

	故选 B。
}

\item
%% number: 4
%% paperName: 2023年海南高考第 4 题 3 分
%% typeId: 1
%% type: 单选题
%% chapter: 机械波
%% point: 波与振动图象
%% method: 无
%% score: 3
%% degree: 550
%% duplicateId: 0
%% body:
下面两图分别是一列机械波在传播方向上相距 $6 \Um$ 的两个质点 P、Q 的振动图像，下列说法正确的是\xzanswer{C}

\twopicture[showlabel=false, wA=5.00cm, wB=5.00cm]{figs/04a.png}{figs/04b.png}

\fourchoices[answer=C]
{该波的周期是 $5 \Us$}
{该波的波速是 $3 \Ums$}
{4 $\Us$ 时 P 质点向上振动}
{4 $\Us$ 时 Q 质点向上振动}

%% answer: C
%% memo:
\memoanswer{
	A．由振动图像可看出该波的周期是 $4 \Us$，A 错误；

	B．由于 Q、P 两个质点振动反相，则可知两者间距离等于 $(n + \frac{1}{2})\lambda = 6$，$n = 0,\ 1,\ 2,\ \cdots$

	根据 $v = \frac{\lambda}{T} = \frac{3}{2n + 1}$，$n = 0,\ 1,\ 2,\ \cdots$。B 错误；

	C．由 P 质点的振动图像可看出，在 $4 \Us$ 时 P 质点在平衡位置向上振动，C 正确；

	D．由 Q 质点的振动图像可看出，在 $4 \Us$ 时 Q 质点在平衡位置向下振动，D 错误。

	故选 C。
}

\item
%% number: 5
%% paperName: 2023年海南高考第 5 题 3 分
%% typeId: 1
%% type: 单选题
%% chapter: 气体性质
%% point: 分子力
%% method: 无
%% score: 3
%% degree: 750
%% duplicateId: 0
%% body:
下列关于分子力和分子势能的说法正确的是\xzanswer{C}

\onepicture[w=6.00cm]{\includesvg{figs/05}}

\fourchoices[answer=C]
{分子间距离大于 $r_{0}$ 时，分子间表现为斥力}
{分子从无限远靠近到距离 $r_{0}$ 处过程中分子势能变大}
{分子势能在 $r_{0}$ 处最小}
{分子间距离小于 $r_{0}$ 且减小时，分子势能在减小}

%% answer: C
%% memo:
\memoanswer{
	分子间距离大于 $r_{0}$，分子间表现为引力，分子从无限远靠近到距离 $r_{0}$ 处过程中，引力做正功，分子势能减小，则在 $r_{0}$ 处分子势能最小；继续减小距离，分子间表现为斥力，分子力做负功，分子势能增大。

	故选 C。
}

\item
%% number: 6
%% paperName: 2023年海南高考第 6 题 3 分
%% typeId: 1
%% type: 单选题
%% chapter: 电磁感应
%% point: 楞次定律推广
%% method: 无
%% score: 3
%% degree: 650
%% duplicateId: 0
%% body:
汽车测速利用了电磁感应现象，汽车可简化为一个矩形线圈 abcd，埋在地下的线圈分别为 1、2，通上顺时针（俯视）方向电流，当汽车经过线圈时\xzanswer{C}

\onepicture[w=5.50cm]{figs/06.png}

\fourchoices[answer=C]
{线圈 1、2 产生的磁场方向竖直向上}
{汽车进入线圈 1 过程产生感应电流方向为 abcd}
{汽车离开线圈 1 过程产生感应电流方向为 abcd}
{汽车进入线圈 2 过程受到的安培力方向与速度方向相同}

%% answer: C
%% memo:
\memoanswer{
	A．由题知，埋在地下的线圈 1、2 通顺时针（俯视）方向的电流，则根据右手定则，可知线圈 1、2 产生的磁场方向竖直向下，A 错误；

	B．汽车进入线圈 1 过程中，磁通量增大，根据楞次定律可知产生感应电流方向为 adcb（逆时针），B 错误；

	C．汽车离开线圈 1 过程中，磁通量减小，根据楞次定律可知产生感应电流方向为 abcd（顺时针），C 正确；

	D．汽车进入线圈 2 过程中，磁通量增大，根据楞次定律可知产生感应电流方向为 adcb（逆时针），再根据左手定则，可知汽车受到的安培力方向与速度方向相反，D 错误。

	故选 C。
}

\item
%% number: 7
%% paperName: 2023年海南高考第 7 题 3 分
%% typeId: 1
%% type: 单选题
%% chapter: 电路
%% point: 电容
%% method: 无
%% score: 3
%% degree: 550
%% duplicateId: 0
%% body:
如图所示电路，已知电源电动势为 $E$，内阻不计，电容器电容为 $C$，闭合开关 K，待电路稳定后，电容器上电荷量为\xzanswer{C}

\onepicture[w=4.50cm]{figs/07.png}

\fourchoices[answer=C]
{$CE$}
{$\frac{1}{2} CE$}
{$\frac{2}{5} CE$}
{$\frac{3}{5} CE$}

%% answer: C
%% memo:
\memoanswer{
	电路稳定后，由于电源内阻不计，则整个回路可看成 $3R$、$2R$ 的串联部分与 $R$、$4R$ 的串联部分并联，若取电源负极为零电势点，则电容器上极板的电势为

	$\varphi_{\text{上}} = \frac{E}{5R}\cdot 2R = \frac{2E}{5}$

	电容器下极板的电势为

	$\varphi_{\text{下}} = \frac{E}{5R}\cdot 4R = \frac{4E}{5}$

	则电容两端的电压

	$U_{\text{下上}} = \frac{2E}{5}$

	则电容器上的电荷量为

	$Q = CU_{\text{上下}} = \frac{2}{5} CE$

	故选 C。
}

\item
%% number: 8
%% paperName: 2023年海南高考第 8 题 3 分
%% typeId: 1
%% type: 单选题
%% chapter: 电场
%% point: 带电粒子的平衡
%% method: 无
%% score: 3
%% degree: 450
%% duplicateId: 0
%% body:
如图所示，一光滑绝缘轨道水平放置，直径上有 A、B 两点，$AO = 2 \Ucm$，$OB = 4 \Ucm$，在 AB 固定两个带电量分别为 $Q_{1}$、$Q_{2}$ 的正电荷，现有一个带正电小球静置于轨道内侧 P 点（小球可视为点电荷），已知 $AP:BP = n:1$，试求 $Q_{1}:Q_{2}$ 是多少\xzanswer{C}

\onepicture[w=5.00cm]{figs/08.png}

\fourchoices[answer=C]
{$2n^{2}:1$}
{$4n^{2}:1$}
{$2n^{3}:1$}
{$4n^{3}:1$}

%% answer: C
%% memo:
\memoanswer{
	对小球受力分析如图所示。由正弦定理有

	$\frac{F_{A}}{\sin\angle CPH} = \frac{F_{B}}{\sin\angle CHP}$

	其中

	$\angle CPH = \angle OPB,\ \angle CHP = \angle HPD = \angle APO$

	其中 $\triangle APO$ 中

	$\frac{AP}{\sin(\pi - \angle POB)} = \frac{AO}{\sin\angle APO}$

	同理有

	$\frac{BP}{\sin\angle POB} = \frac{BO}{\sin\angle BPO}$

	其中

	\onepicture[w=4.50cm]{figs/08_an.png}

	$F_{A} = k\frac{Q_{1}q}{AP^{2}},\quad F_{B} = k\frac{Q_{2}q}{BP^{2}}$

	联立有

	$Q_{1}:Q_{2} = 2n^{3}:1$

	故选 C。
}

\item
%% number: 9
%% paperName: 2023年海南高考第 9 题 4 分
%% typeId: 2
%% type: 多选题
%% chapter: 曲线运动
%% point: 人造地球卫星
%% method: 无
%% score: 4
%% degree: 650
%% duplicateId: 0
%% body:
如图所示，1、2 轨道分别是天宫二号飞船在变轨前后的轨道，下列说法正确的是\xzanswer{ACD}

\onepicture[w=4.00cm]{\includesvg{figs/09}}

\fourchoices[answer=ACD]
{飞船从 1 轨道变到 2 轨道要点火加速}
{飞船在 1 轨道周期大于 2 轨道周期}
{飞船在 1 轨道速度大于 2 轨道}
{飞船在 1 轨道加速度大于 2 轨道}

%% answer: ACD
%% memo:
\memoanswer{
	A．飞船从较低的轨道 1 进入较高的轨道 2 要进行加速做离心运动才能完成，选项 A 正确；

	BCD．根据

	$G\frac{Mm}{r^{2}} = m\frac{v^{2}}{r} = m\frac{4\pi^{2}}{T^{2}}r = ma$

	可得 $a = \frac{GM}{r^{2}}$，$v = \sqrt{\frac{GM}{r}}$，$T = 2\pi\sqrt{\frac{r^{3}}{GM}}$

	可知飞船在轨道 1 的周期小于在轨道 2 的周期，在轨道 1 的速度大于在轨道 2 的速度，在轨道 1 的加速度大于在轨道 2 的加速度，故选项 B 错误，CD 正确。

	故选 ACD。
}

\item
%% number: 10
%% paperName: 2023年海南高考第 10 题 4 分
%% typeId: 2
%% type: 多选题
%% chapter: 光学
%% point: 光子说
%% method: 无
%% score: 4
%% degree: 650
%% duplicateId: 0
%% body:
已知一个激光发射器功率为 $P$，发射波长为 $\lambda$ 的光，光速为 $c$，普朗克常量为 $h$，则\xzanswer{AC}

\fourchoices[answer=AC]
{光的频率为 $\frac{c}{\lambda}$}
{光子的能量为 $\frac{h}{\lambda}$}
{光子的动量为 $\frac{h}{\lambda}$}
{在时间 $t$ 内激光器发射的光子数为 $\frac{Ptc}{h\lambda}$}

%% answer: AC
%% memo:
\memoanswer{
	A．光的频率 $\nu = \frac{c}{\lambda}$。选项 A 正确；

	B．光子的能量 $E = h\nu = \frac{hc}{\lambda}$。选项 B 错误；

	C．光子的动量 $p = \frac{h}{\lambda}$。选项 C 正确；

	D．在时间 $t$ 内激光器发射的光子数 $n = \frac{Pt}{E} = \frac{Pt\lambda}{hc}$。选项 D 错误。

	故选 AC。
}

\item
%% number: 11
%% paperName: 2023年海南高考第 11 题 4 分
%% typeId: 2
%% type: 多选题
%% chapter: 交流电
%% point: 交流电
%% method: 无
%% score: 4
%% degree: 650
%% duplicateId: 0
%% body:
下图是工厂利用 $u = 220\sqrt{2}\sin 100\pi t \UV$ 的交流电给 $36 \UV$ 照明灯供电的电路，变压器原线圈匝数为 1 100 匝，下列说法正确的是\xzanswer{BC}

\onepicture[w=6.50cm]{figs/11.png}

\fourchoices[answer=BC]
{电源电压有效值为 $220\sqrt{2} \UV$}
{交变电流的周期为 $0.02 \Us$}
{副线圈匝数为 180 匝}
{副线圈匝数为 240 匝}

%% answer: BC
%% memo:
\memoanswer{
	A．电源电压的有效值 $U = \frac{220\sqrt{2}}{\sqrt{2}} \UV = 220 \UV$。选项 A 错误；

	B．交流电的周期 $T = \frac{2\pi}{\omega} = \frac{2\pi}{100\pi} \Us = 0.02 \Us$。选项 B 正确；

	CD．根据 $\frac{n_{1}}{n_{2}} = \frac{U_{1}}{U_{2}}$ 可得副线圈匝数 $n_{2} = \frac{U_{2}}{U_{1}} n_{1} = \frac{36}{220} \times 1100 = 180$ 匝。选项 C 正确，D 错误。

	故选 BC。
}

\item
%% number: 12
%% paperName: 2023年海南高考第 12 题 4 分
%% typeId: 2
%% type: 多选题
%% chapter: 电场
%% point: 电势能电势差电场力做功
%% method: 无
%% score: 4
%% degree: 550
%% duplicateId: 0
%% body:
如图所示，正三角形三个顶点固定三个等量电荷，其中 A、B 带正电，C 带负电，O、M、N 为 AB 边的四等分点，下列说法正确的是\xzanswer{BC}

\onepicture[w=5.50cm]{figs/12.png}

\fourchoices[answer=BC]
{M、N 两点电场强度相同}
{M、N 两点电势相同}
{负电荷在 M 点电势能比在 O 点时要小}
{负电荷在 N 点电势能比在 O 点时要大}

%% answer: BC
%% memo:
\memoanswer{
	A．根据场强叠加以及对称性可知，MN 两点的场强大小相同，但是方向不同，选项 A 错误；

	B．因在 AB 处的正电荷在 MN 两点的合电势相等，在 C 点的负电荷在 MN 两点的电势也相等，则 MN 两点电势相等，选项 B 正确；

	CD．因负电荷从 M 到 O，因 AB 两电荷的合力对负电荷的库仑力从 O 指向 M，则该力对负电荷做负功，C 点的负电荷也对该负电荷做负功，可知三个电荷对该负电荷的合力对其做负功，则该负电荷的电势能增加，即负电荷在 M 点的电势能比在 O 点小；同理可知负电荷在 N 点的电势能比在 O 点小。选项 C 正确，D 错误。

	故选 BC。
}

\item
%% number: 13
%% paperName: 2023年海南高考第 13 题 4 分
%% typeId: 2
%% type: 多选题
%% chapter: 磁场
%% point: 带电粒子在组合场中的运动
%% method: 无
%% score: 4
%% degree: 450
%% duplicateId: 0
%% body:
如图所示，质量为 $m$，带电量为 $+q$ 的点电荷，从原点以初速度 $v_{0}$ 射入第一象限内的电磁场区域，在 $0 < y < y_{0}$，$0 < x < x_{0}$（$x_{0}$、$y_{0}$ 为已知）区域内有竖直向上的匀强电场，在 $x > x_{0}$ 区域内有垂直纸面向里的匀强磁场，控制电场强度（$E$ 值有多种可能），可让粒子从 NP 射入磁场后偏转打到接收器 MN 上，则\xzanswer{AD}

\onepicture[w=6.50cm]{figs/13.png}

\fourchoices[answer=AD]
{粒子从 NP 中点射入磁场，电场强度满足 $E = \frac{y_{0}mv_{0}^{2}}{qx_{0}^{2}}$}
{粒子从 NP 中点射入磁场时速度为 $v_{0}\sqrt{\frac{x_{0}^{2} + y_{0}^{2}}{y_{0}^{2}}}$}
{粒子在磁场中做圆周运动的圆心到 NM 的距离为 $\frac{mv_{0}}{qB}$}
{粒子在磁场中运动的圆周半径最大值是 $\frac{mv_{0}}{qB}\sqrt{\frac{x_{0}^{2} + 4y_{0}^{2}}{x_{0}^{2}}}$}

%% answer: AD
%% memo:
\memoanswer{
	A．若粒子打到 PN 中点，则

	$x_{0} = v_{0}t_{1}$

	$\frac{1}{2}y_{0} = \frac{1}{2}\cdot\frac{qE}{m}t^{2}$

	解得 $E = \frac{y_{0}mv_{0}^{2}}{qx_{0}^{2}}$。选项 A 正确；

	B．粒子从 PN 中点射出时，则 $\frac{y_{0}}{2} = \frac{v_{y}}{2} t_{1}$

	速度 $v_{1} = \sqrt{v_{0}^{2} + v_{y}^{2}} = v_{0}\sqrt{\frac{x_{0}^{2} + y_{0}^{2}}{x_{0}^{2}}}$。选项 B 错误；

	C．粒子从电场中射出时的速度方向与竖直方向夹角为 $\theta$，则

	$\tan\theta = \frac{v_{0}}{v_{y}} = \frac{v_{0}}{\frac{qE}{m}\cdot\frac{x_{0}}{v_{0}}} = \frac{mv_{0}^{2}}{qEx_{0}}$

	粒子从电场中射出时的速度

	$v = \frac{v_{0}}{\sin\theta}$

	粒子进入磁场后做匀速圆周运动，则

	$qvB = m\frac{v^{2}}{r}$

	则粒子进入磁场后做圆周运动的圆心到 MN 的距离为 $d = r\cos\theta$

	解得 $d = \frac{Ex_{0}}{Bv_{0}}$

	\onepicture[w=6.50cm]{figs/13_an.jpg}

	选项 C 错误；

	D．当粒子在磁场中运动有最大运动半径时，进入磁场的速度最大，则此时粒子从 N 点进入磁场，此时竖直最大速度

	$v_{ym} = \frac{2y_{0}}{t}$

	$x_{0} = v_{0}t$

	出离电场的最大速度 $v_{m} = \sqrt{v_{0}^{2} + v_{ym}^{2}} = v_{0}\sqrt{\frac{x_{0}^{2} + 4y_{0}^{2}}{x_{0}^{2}}}$

	则由 $qvB = m\frac{v^{2}}{r}$

	可得最大半径 $r_{m} = \frac{mv_{m}}{qB} = \frac{mv_{0}}{qB}\sqrt{\frac{x_{0}^{2} + 4y_{0}^{2}}{x_{0}^{2}}}$

	选项 D 正确；

	故选 AD。
}

\item
%% number: 14
%% paperName: 2023年海南高考第 14 题 6 分
%% typeId: 4
%% type: 实验题
%% chapter: 光学
%% point: 测定玻璃折射率
%% method: 无
%% score: 6
%% degree: 550
%% duplicateId: 0
%% body:
用激光测玻璃砖折射率的实验中，玻璃砖与屏 P 平行放置，从另一侧用激光笔以一定角度照射，此时在屏上的 $S_{1}$ 处有激光点，移走玻璃砖，光点移到 $S_{2}$ 处，回答下列问题：
\begin{enumerate}
	\item
	请画出激光束经玻璃折射后完整的光路图\tkanswer{（如图）}；
	\item
	已经测出 $AB = l_{1}$，$OA = l_{2}$，$S_{1}S_{2} = l_{3}$，则折射率 $n = $\tkanswer{$\frac{l_{1}\sqrt{(l_{1} - l_{3})^{2} + l_{2}^{2}}}{(l_{1} - l_{3})\sqrt{l_{1}^{2} + l_{2}^{2}}}$}（用 $l_{1}$、$l_{2}$、$l_{3}$ 表示）；
	\item
	若改用宽 ab 更小的玻璃砖做实验，则 $S_{1}S_{2}$ 间的距离会\tkanswer{变小}（填“变大”，“变小”或“不变”）。
\end{enumerate}
\onepicture[w=6.00cm, align=right]{figs/14.png}
%% answer: 
\jdanswer{
\begin{enumerate}
	\item
	光路图如图所示。
	\item
	$\frac{l_{1}\sqrt{(l_{1} - l_{3})^{2} + l_{2}^{2}}}{(l_{1} - l_{3})\sqrt{l_{1}^{2} + l_{2}^{2}}}$
	\item
	变小
\end{enumerate}
}
%% memo:
\memoanswer{
	（1）根据题意画出光路图如下图所示

	\onepicture[w=5.00cm]{figs/14_an1.png}

	（2）设光线入射角为 $\theta$、折射角为 $\alpha$，则在 C 点根据折射定律有 $n\sin\theta = \sin\alpha$

	由于射入玻璃砖的入射角是射出玻璃砖的折射角，则 $S_{1}S_{2} = CB$

	根据几何关系可知 $\sin\theta = \frac{l_{1} - l_{3}}{\sqrt{(l_{1} - l_{3})^{2} + l_{2}^{2}}}$，$\sin\alpha = \frac{l_{1}}{\sqrt{l_{1}^{2} + l_{2}^{2}}}$

	联立解得 $n = \frac{l_{1}\sqrt{(l_{1} - l_{3})^{2} + l_{2}^{2}}}{(l_{1} - l_{3})\sqrt{l_{1}^{2} + l_{2}^{2}}}$

	（3）若改用宽 ab 更小的玻璃砖做实验，则画出光路图如下

	\onepicture[w=5.00cm]{figs/14_an2.png}

	可看出 $S_{1}S_{2}$ 间的距离变小。
}

\item
%% number: 15
%% paperName: 2023年海南高考第 15 题 14 分
%% typeId: 4
%% type: 实验题
%% chapter: 电路
%% point: 电表改装
%% method: 无
%% score: 14
%% degree: 450
%% duplicateId: 0
%% body:
用如图~\ref{2023海南15a}所示的电路测量一个量程为 $100 \UuA$，内阻约为 $2\,000 \UO$ 的微安表头的内阻，所用电源的电动势约为 $12 \UV$，有两个电阻箱可选，$R_{1}$（$0 \sim 9999.9 \UO$），$R_{2}$（$99999.9 \UO$）。

\begin{enumerate}
	\item
	$R_{M}$ 应选\tkanswer{$R_{1}$}，$R_{N}$ 应选\tkanswer{$R_{2}$}；
	\item
	根据电路图，请把实物连线补充完整\tkanswer{（如图）}；
	\item
	下列操作顺序合理排列是\tkanswer{①③②④}：
	\begin{steps}
		\item
		将变阻器滑动头 P 移至最左端，将 $R_{N}$ 调至最大值；
		\item
		闭合开关 $S_{2}$，调节 $R_{M}$，使微安表半偏，并读出 $R_{M}$ 阻值；
		\item
		断开 $S_{2}$，闭合 $S_{1}$，调节滑动头 P 至某位置再调节 $R_{N}$ 使表头满偏；
		\item
		断开 $S_{1}$、$S_{2}$，拆除导线，整理好器材
	\end{steps}
	\item
	如图~\ref{2023海南15c}是 $R_{M}$ 调节后面板，则待测表头的内阻为\tkanswer{$1\,998.0 \UO$}，该测量值\tkanswer{小于}（填“大于”、“小于”、“等于”）真实值。
	\item
	将该微安表改装成量程为 $2 \UV$ 的电压表后，某次测量指针指在图~\ref{2023海南15d}所示位置，则待测电压为\tkanswer{1.28} $\UV$（保留 3 位有效数字）。
	\item
	某次半偏法测量表头内阻的实验中，$S_{2}$ 断开，电表满偏时读出 $R_{N}$ 值，在滑动头 P 不变，$S_{2}$ 闭合后调节电阻箱 $R_{M}$，使电表半偏时读出 $R_{M}$，若认为 OP 间电压不变，则微安表内阻为\tkanswer{$\frac{R_{N}R_{M}}{R_{N} - R_{M}}$}（用 $R_{M}$、$R_{N}$ 表示）
\end{enumerate}
\fourpicture[labelA=2023海南15a, labelB=2023海南15b, labelC=2023海南15c, labelD=2023海南15d, hA=3.60cm, hB=5.00cm, hC=4.50cm, hD=2.60cm, align=right]
{figs/15a.png}
{figs/15b.png}
{figs/15c.png}
{figs/15d.png}
%% answer: 
\jdanswer{
\begin{enumerate}
	\item
	$R_{1}$，$R_{2}$
	\item
	实物连线如图所示。
	\item
	①③②④
	\item
	$1\,998.0 \UO$，小于
	\item
	1.28
	\item
	$\frac{R_{N}R_{M}}{R_{N} - R_{M}}$
\end{enumerate}
}
%% memo:
\memoanswer{
	（1）根据半偏法的测量原理可知，$R_{M}$ 与 $R_{1}$ 相当，当闭合 $S_{2}$ 之后，变阻器上方的电流应基本不变，就需要 $R_{N}$ 较大，对下方分压电路影响甚微。故 $R_{M}$ 应选 $R_{1}$，$R_{N}$ 应选 $R_{2}$。

	（2）根据电路图连接实物图有

	\onepicture[w=6.50cm]{figs/15_an1.png}

	（3）根据半偏法的实验步骤应为

	① 将变阻器滑动头 P 移至最左端，将 $R_{N}$ 调至最大值；

	③ 断开 $S_{2}$，闭合 $S_{1}$，调节滑动头 P 至某位置再调节 $R_{N}$ 使表头满偏；

	② 闭合开关 $S_{2}$，调节 $R_{M}$，使微安表半偏，并读出 $R_{M}$ 阻值；

	④ 断开 $S_{1}$、$S_{2}$，拆除导线，整理好器材。

	（4）根据 $R_{M}$ 调节后面板读数为 $1\,998.0 \UO$。

	当闭合 $S_{2}$ 后，原电路可看成如下电路

	\onepicture[w=7.00cm]{figs/15_an2.png}

	闭合 $S_{2}$ 后，相当于 $R_{M}$ 由无穷大变成有限值，变小了，则流过 $R_{N}$ 的电流大于原来的电流，则流过 $R_{M}$ 的电流大于 $\frac{I_{\text{原}}}{2}$，故待测表头的内阻的测量值小于真实值。

	（5）将该微安表改装成量程为 $2 \UV$ 的电压表，则需要串联一个电阻 $R_{0}$，则有

	$U = I_{g}(R_{g} + R_{0})$

	此时的电压读数有 $U' = I'(R_{g} + R_{0})$

	其中 $U = 2 \UV$，$I_{g} = 100 \UuA$，$I' = 64 \UuA$

	联立解得 $U' = 1.28 \UV$

	（6）根据题意 OP 间电压不变，可得

	$I(R_{A} + R_{N}) = \left(\frac{I}{2} + \frac{\frac{I}{2}R_{A}}{R_{M}}\right)R_{N} + \frac{I}{2}R_{A}$

	解得 $R_{A} = \frac{R_{N}R_{M}}{R_{N} - R_{M}}$
}

\item
%% number: 16
%% paperName: 2023年海南高考第 16 题 8 分
%% typeId: 6
%% type: 计算题
%% chapter: 气体性质
%% point: 气体多过程
%% method: 无
%% score: 8
%% degree: 750
%% duplicateId: 0
%% body:
某饮料瓶内密封一定质量理想气体，$t = 27\UdC$ 时，压强 $p = 1.050 \times 10^{5} \UPa$。
\begin{enumerate}
	\item
	$t' = 37\UdC$ 时，气压是多大？
	\item
	保持温度不变，挤压气体，使之压强与（1）时相同时，气体体积为原来的多少倍？
\end{enumerate}
\onepicture[h=4.00cm, align=right]{figs/16.png}
%% answer: 
\jdanswer{
\begin{enumerate}
	\item
	$1.085 \times 10^{5} \UPa$
	\item
	0.97
\end{enumerate}
}
%% memo:
\memoanswer{
	（1）瓶内气体的始末状态的热力学温度分别为

	$T = (27 + 273) \UK = 300 \UK$，$T' = (37 + 273) \UK = 310 \UK$

	温度变化过程中体积不变，故由查理定律有 $\frac{p}{T} = \frac{p'}{T'}$

	解得 $p' = 1.085 \times 10^{5} \UPa$

	（2）保持温度不变，挤压气体，等温变化过程，由玻意耳定律有

	$pV = p'V'$

	解得 $V' \approx 0.97V$
}

\item
%% number: 17
%% paperName: 2023年海南高考第 17 题 12 分
%% typeId: 6
%% type: 计算题
%% chapter: 磁场
%% point: 安培力
%% method: 无
%% score: 12
%% degree: 550
%% duplicateId: 0
%% body:
如图所示，U 形金属杆上边长为 $L = 15 \Ucm$，质量为 $m = 1 \times 10^{-3} \Ukg$，下端插入导电液体中，导电液体连接电源，金属杆所在空间有垂直纸面向里 $B = 8 \times 10^{-2} \UT$ 的匀强磁场。
\begin{enumerate}
	\item
	若插入导电液体部分深 $h = 2.5 \Ucm$，闭合电键后，金属杆飞起后，其下端离液面高度 $H = 10 \Ucm$，设杆中电流不变，求金属杆离开液面时的速度大小和金属杆中的电流有多大；（$g$ 取 $10 \Umsq$）
	\item
	若金属杆下端刚与导电液体接触，改变电动势的大小，通电后金属杆跳起高度 $H' = 5 \Ucm$，通电时间 $t' = 0.002 \Us$，求通过金属杆截面的电荷量。
\end{enumerate}
\onepicture[w=7.00cm, align=right]{figs/17.png}
%% answer: 
\jdanswer{
\begin{enumerate}
	\item
	$\sqrt{2} \Ums$，4.17 $\UA$
	\item
	0.085 $\UC$
\end{enumerate}
}
%% memo:
\memoanswer{
	（1）对金属杆，跳起的高度为 $H$，竖直上抛运动由运动学关系式 $v^{2} = 2gH$，解得 $v = \sqrt{2gH} = \sqrt{2} \Ums$

	通电过程金属杆受到的安培力大小为 $F_{A} = BIL$

	由动能定理得 $BILh - mg(H + h) = 0$

	解得 $I = 4.17 \UA$

	（2）对金属杆，通电时间 $t' = 0.002 \Us$，由动量定理有 $(BI'L - mg)t' = mv' - 0$

	由运动学公式 $v'^{2} = 2gH'$

	通过金属杆截面的电荷量 $q = I't'$

	联立解得 $q = 0.085 \UC$
}

\item
%% number: 18
%% paperName: 2023年海南高考第 18 题 16 分
%% typeId: 6
%% type: 计算题
%% chapter: 运动和力
%% point: 牛顿定律的应用
%% method: 无
%% score: 16
%% degree: 450
%% duplicateId: 0
%% body:
如图所示，有一固定的光滑 $\frac{1}{4}$ 圆弧轨道，半径 $R = 0.2 \Um$，一质量为 $m_{B} = 1 \Ukg$ 的小滑块 B 从轨道顶端滑下，在其冲上长木板 C 左端时，给木板一个与小滑块相同的初速度，已知 $m_{C} = 3 \Ukg$，B、C 间动摩擦因数 $\mu_{1} = 0.2$，C 与地面间的动摩擦因数 $\mu_{2} = 0.8$，C 右端有一个挡板，C 长为 $L$。求：
\begin{enumerate}
	\item
	B 滑到 A 的底端时对 A 的压力是多大？
	\item
	若 B 未与 C 右端挡板碰撞，当 B 与地面保持相对静止时，B、C 间因摩擦产生的热量是多少？
	\item
	在 $0.16 \Um < L < 0.8 \Um$ 时，B 与 C 右端挡板发生碰撞，且碰后粘在一起，求 B 从滑上 C 到最终停止所用的时间。
\end{enumerate}
\onepicture[w=8.00cm, align=right]{figs/18.png}
%% answer: 
\jdanswer{
\begin{enumerate}
	\item
	30 $\UN$
	\item
	1.6 $\UJ$
	\item
	$\left(1 - \frac{15\sqrt{0.8 - L}}{16}\right) \Us$
\end{enumerate}
}
%% memo:
\memoanswer{
	（1）滑块下滑到轨道底部，有

	$m_{B}gR = \frac{1}{2}m_{B}v_{0}^{2} - 0$

	解得 $v_{0} = 2 \Ums$

	在底部，根据牛顿第二定律

	$F_{N} - m_{B}g = m_{B}\frac{v_{0}^{2}}{R}$

	解得 $F_{N} = 30 \UN$

	由牛顿第三定律可知 B 对 A 的压力是 $30 \UN$。

	（2）当 B 滑上 C 后，对 B 分析，受摩擦力向左，根据牛顿第二定律得

	$f_{B} = \mu_{1}m_{B}g = m_{B}a_{B}$

	解得加速度向左为 $a_{1} = 2 \Umsq$

	对 C 分析，受 B 向右的摩擦力 $\mu_{1}m_{B}g$ 和地面向左的摩擦力

	$f_{\text{地}C} = \mu_{2}(m_{B} + m_{C})g$

	根据牛顿第二定律

	$f_{\text{地}C} - f_{BC} = \mu_{2}(m_{B} + m_{C})g - \mu_{1}m_{B}g = m_{C}a_{2}$

	解得其加速度向左为 $a_{2} = 10 \Umsq$

	由运动学位移与速度关系公式 $v^{2} - v_{0}^{2} = 2ax$，得 B 向右运动的距离 $x_{1} = \frac{v_{0}^{2}}{2a_{1}} = 1 \Um$

	C 向右运动距离 $x_{2} = \frac{v_{0}^{2}}{2a_{2}} = 0.2 \Um$

	由功能关系可知，B、C 间摩擦产生的热量 $Q = \mu_{1}m_{B}g(x_{1} - x_{2})$

	可得 $Q = 1.6 \UJ$

	（3）由上问可知，若 B 还未与 C 上挡板碰撞，C 先停下，用时为 $t_{1}$，有 $t_{1} = \frac{v_{0}}{a_{2}}$

	解得 $t_{1} = 0.2 \Us$

	B 的位移为 $x_{B1} = v_{0}t_{1} - \frac{1}{2}a_{1}t_{1}^{2} = 0.36 \Um$

	则此刻的相对位移为 $x_{\text{相}} = x_{B1} - x_{2} = 0.16 \Um$

	此时 $v_{B1} = v_{0} - a_{1}t_{1} = 1.6 \Ums$

	由 $L > 0.16 \Um$，一定是 C 停下之后，B 才与 C 上挡板碰撞。设再经 $t_{2}$ 时间 B 与 C 挡板碰撞，有

	$L - 0.16 = 1.6t_{2} - \frac{1}{2}a_{1}t_{2}^{2}$

	解得 $t_{2} = 0.8 - \sqrt{0.8 - L}$

	碰撞时 B 速度为 $v_{B2} = v_{B1} - a_{1}t_{2} = 2\sqrt{0.8 - L}$

	碰撞时由动量守恒可得 $m_{B}v_{B2} = (m_{B} + m_{C})v$

	解得碰撞后 B、C 速度为

	$v = \frac{\sqrt{0.8 - L}}{2}$

	之后二者一起减速，根据牛顿第二定律得

	$a_{3} = \frac{\mu_{2}(m_{B} + m_{C})g}{m_{B} + m_{C}} = 8 \Umsq$

	再经 $t_{3}$ 后停下，则有

	$t_{3} = \frac{v}{a_{3}} = \frac{\sqrt{0.8 - L}}{16}$

	故 B 从滑上 C 到最终停止所用的时间总时间 $t = t_{1} + t_{2} + t_{3} = \left(1 - \frac{15\sqrt{0.8 - L}}{16}\right) \Us$
}

\end{enumerate}

\end{document}
